\section{Digital Music as a Practical Software Domain}\label{sec:digital-music-domain}

Recorded music is now a large digital ecosystem rather than a niche computing topic.
IFPI reported that global recorded music revenues reached US\$31.7 billion in 2025, with streaming representing 69.6\%
of total recorded music revenue and paid subscription streaming representing 52.4\% of the total
market~\cite{IFPI_2026}.
Spotify, reporting from the perspective of a single streaming platform, stated that it paid more than US\$11 billion to
the music industry in 2025~\cite{Spotify_2026_Payouts}.
BandLab's public product material describes a community of 100 million music creators using its platform
ecosystem~\cite{BandLab_2026}.
These figures should not be read as a direct measurement of demand for a programming language.
They do, however, show that digital music creation, distribution, and consumption are large enough for specialized
creative software to be a relevant engineering topic.

\begin{figure}[htbp]
  \centering
  \motivoscalefig{%
    \begin{tikzpicture}[
        fact/.style={rounded corners=4pt, draw=white, line width=0.9pt, minimum width=3.7cm, minimum height=2.1cm,
        align=center, font=\sffamily, text=white},
        caption/.style={font=\scriptsize\sffamily, text=motivoInk, align=center}
      ]
      \node[fact, fill=motivoBlue] (revenue) at (0, 0) {\Large \$31.7B\\[0.15cm]\normalsize global recorded\\music
      revenue\\in 2025};
      \node[fact, fill=motivoCyan] (streaming) at (4.5, 0) {\Large 69.6\%\\[0.15cm]\normalsize revenue share\\from
      streaming};
      \node[fact, fill=motivoViolet] (creators) at (9.0, 0) {\Large 100M\\[0.15cm]\normalsize BandLab music\\creators};

      \node[caption, below=0.3cm of revenue] {IFPI Global Music\\Report 2026};
      \node[caption, below=0.3cm of streaming] {IFPI format\\breakdown};
      \node[caption, below=0.3cm of creators] {BandLab public\\platform material};
    \end{tikzpicture}%
  }
  \caption{Industry-scale signals motivating digital music tooling.
    The numbers are attributed to their reporting organizations and are used as context, not as a direct evaluation of
  Motivo.}\label{fig:industry-context}
\end{figure}

The production side of this ecosystem is heavily tool-based.
Digital audio workstations make detailed editing possible through piano rolls, arrangement timelines, automation lanes,
plug-ins, and MIDI routing.
For many tasks this visual approach is the right interface.
It is direct, immediate, and close to the way a composer hears a piece.
At the same time, repeated symbolic structures can become awkward when they are represented only as copied note blocks.
Changing the duration, transposition, or placement rule of a recurring phrase may require many local edits.
The visible notes remain editable, but the reason they exist is no longer represented as a reusable rule.

\subsection{A Brief History of Musical Programming Languages}\label{subsec:musical-pl-history}

Computer-assisted music composition predates personal digital audio workstations.
Mathews and collaborators developed the \textit{MUSIC} family of programs at Bell Labs from the late 1950s onward,
establishing the idea that musical structures could be described as programs and rendered by a
machine~\cite{Roads_1996}.
Subsequent decades brought specialized environments for sound synthesis, algorithmic composition, and interactive
performance.
Roads surveys this landscape as a continuum from symbolic score generation through real-time digital signal
processing~\cite{Roads_1996}.
Nierhaus organizes algorithmic composition more narrowly as rule- and process-driven generation of musical
material~\cite{Nierhaus_2009}.

Two developments are especially relevant to the present work.
First, the Musical Instrument Digital Interface (MIDI), introduced in 1983, standardized a compact event representation
for note onsets, offsets, and control changes~\cite{Moog_1986}.
MIDI made it practical to exchange symbolic music between instruments and computers, but it encodes performance events
rather than compositional abstractions~\cite{Moore_1988}.
Second, the live-coding movement of the 2000s and 2010s produced languages such as \textit{TidalCycles} and
\textit{Sonic Pi} that treat code as a performance instrument~\cite{McLean_Wiggins_2010,Aaron_2016}.
These systems emphasize real-time reaction, pattern transformation, and integration with synthesis backends.

Across this history, most musical programming languages have been designed for \textit{live performance} or
\textit{audio synthesis}.
They prioritize low-latency scheduling, concurrent processes, and sound generation.
Relatively few systems treat a Standard MIDI File as the primary, deterministic compilation target for offline
structural composition.
Motivo occupies that narrower position: a compiled language whose main artifact is a portable MIDI file, not a running
audio engine.

Algorithmic composition takes the opposite direction from manual piano-roll editing by describing music through formal
procedures.
Motivo adopts only one part of that tradition.
It does not attempt to synthesize audio, model performer expression, or improvise during playback.
It focuses on symbolic generation: a program should produce a MIDI file whose notes can then be inspected, played, or
arranged further in any MIDI-aware environment.
